\begin{textsample}{POS Dim 1 – ma – Score 67.00 – ma/ma\_em\_32.txt}  \label{ex:f1_pos_061}
3. \textbf{Um} olhar mais profundo sobre as áreas de conhecimentos : seus componentes, competências e habilidades 3.1.
Área de Linguagem e suas Tecnologias A comunicação humana é a prática \textbf{que} permite as interações sociais e pode ocorrer tanto de modo formal quanto informal.
Para isso, fazemos uso também da linguagem.
Sendo assim, a linguagem deve ser concebida como \textbf{uma} parte integrante da vida social, pois é considerada o instrumento principal da comunicação informativa.
Para \textbf{que} a linguagem ocorra de forma coerente, fazemos uso dos meios de comunicação, \textbf{que}, na verdade, são instrumentos \textbf{que} possibilitam a transmissão e troca de informações entre duas ou mais pessoas.
Nesse contexto, cabe considerarmos os elementos da comunicação, cada aspecto presente no fluxo comunicativo, apresentados por Roman Jakobson.
Sabemos \textbf{que} os elementos da comunicação propostos por Jakobson estão presentes em qualquer tipo de comunicação e podem, nesse sentido, sofrer alterações para se adequarem ao contexto de prática comunicativa : 1. emissor ou locutor — quem elabora e emite a mensagem ; 2. receptor ou interlocutor — quem recebe a mensagem emitida ; 3. mensagem — conjunto de informações, verbais ou não verbais, transmitidas pelo emissor ; 4. referente ou contexto — situação comunicativa em \textbf{que} se encontram os interlocutores, ou seja, ambiente, espaço, tempo ; 5. canal ou veículo — o meio pelo qual a mensagem é difundida, divulgada ; 6. código — conjunto de signos organizados e utilizados na transmissão da mensagem.
Portanto, para compreensão e uso dos elementos da comunicação e, consequentemente, para garantirmos a realização da comunicação da forma mais coerente possível, inegavelmente, torna -se importante o estudo e aprendizado de competências e habilidades das áreas de conhecimento, neste caso, de forma mais pontual, da área de Linguagem e suas Tecnologias.
Para isso, primeiramente, devemos entender em \textbf{que} palco se encontra esta discussão na BNCC, nos documentos educacionais oficiais do país e como será a proposta de trabalho do estado do Maranhão.
Além disso, dentro da arquitetura curricular dessa área, os conteúdos deverão ser \textbf{embasados} de forma a contribuir para o alcance dos objetivos dos estudantes : continuidade dos estudos, atuação no mundo do trabalho, resolução de demandas complexas da vida cotidiana e exercício pleno da cidadania.
Para isso, cabe mencionarmos \textbf{que} o aprofundamento dos conhecimentos e habilidades da área tem, nos itinerários formativos, ferramentas \textbf{que} ajudam a consolidação desses conhecimentos em diferentes contextos.
Ressaltamos \textbf{que}, por meio dos componentes curriculares na BNCC, na área da linguagem é previsto \textbf{que} : > os estudantes desenvolvam competências e habilidades \textbf{que} lhes possibilitem mobilizar e articular conhecimentos desses componentes, simultaneamente a dimensões socioemocionais, em situações de aprendizagem \textbf{que} lhes sejam significativas e relevantes para sua formação integral ( BRASIL, 2017, p. 481 ).
Podemos iniciar \textbf{explicitando} \textbf{que} a área de Linguagens e suas Tecnologias tem como função primeira a comunicação, por meio de \textbf{um} sistema repleto de signos, símbolos e formas \textbf{que} se apresentam como resultado das relações de comunicação ( verbal, visual, corporal, sonora, icônica, digital, entre outras ) \textbf{que} envolvem o processo cognitivo e criativo necessário à apropriação e compreensão da comunicação humana.
A comunicação se estrutura, na língua portuguesa ( língua materna ), como geradora de significação e, também, integradora na organização das relações humanas no mundo, por meio da aquisição e domínio dos signos da língua escrita e falada.
A língua estrangeira se apresenta como possibilidade de integração, acesso e inclusão inter e intraterritorial, junto a pessoas e culturas diversas.
A arte, como área do conhecimento, por meio das quatro linguagens ( teatro, artes visuais, dança e música ), como expressão criadora e geradora de significação de \textbf{uma} linguagem \textbf{que} possibilita ao estudante \textbf{uma} formação crítica em decorrência da ação mediada do “ \textbf{ver} ”, do “ fazer ”, do “ ouvir ” e do “ sentir ”.
A educação física, por sua vez, propõe o domínio do corpo como forma de expressão e de qualidade de vida.
A educação física se propõe a tematizar as práticas corporais em suas diversas manifestações, explorando movimentos e gestos de diferentes culturas e sociedades, possibilitando novas experiências aos estudantes.
Sendo assim, podemos \textbf{dizer} \textbf{que} o conjunto de componentes curriculares reunidos na área de Linguagens e suas Tecnologias – somado às informações adquiridas das demais áreas do conhecimento – permite ao estudante, quando mediado adequadamente pelo docente, \textbf{uma} formação completa, conforme \textbf{preconiza} a BNCC ( BRASIL, 2017 ) quando afirma \textbf{que} “ por meio dessas práticas, as pessoas interagem consigo mesmas e com os outros, constituindo -se como \textbf{sujeitos} sociais.
Nessas interações, estão \textbf{imbricados} conhecimentos, atitudes e valores culturais, morais e éticos ” ( p. 62 ).
É importante entendermos, também, \textbf{que} a unificação de componentes em áreas foi pensada e planejada como estratégia para organizar os conteúdos e interligar as \textbf{disciplinas} para \textbf{uma} atuação conjunta, possibilitando a \textbf{transdisciplinaridade} ( termo atribuído a Jean Piaget na década de 1970 ).
Neste formato, os professores podem contextualizar os conteúdos e abordá -los concomitantemente, com o intuito de propiciar ao estudante o desenvolvimento de habilidades.
Para isso, os conteúdos devem ser pensados e \textbf{ancorados} nos documentos oficiais, formalizando suas aplicações em âmbito nacional.
Portanto, a área de Linguagem no ensino médio, considerando as aprendizagens adquiridas no ensino fundamental, tem o desafio de aprofundar possibilidades de uso da língua verbal ( oral ou \textbf{visual-motora}, como Libras e escrita ), considerando a diversidade étnica dos diversos \textbf{sujeitos} ( indígenas, quilombolas, \textbf{ribeirinhos}, campesinos, ciganos, imigrantes, entre outros ), a pessoa com deficiência, transtorno do espectro autista, altas habilidades e superdotação, de modo \textbf{que} os conhecimentos linguísticos contribuam para ampliar a compreensão e comunicação territorial e global, por meio das línguas materna e estrangeiras ( inglês e espanhol ), respeitando as diversas culturas no mundo e, ainda, os diversos dialetos falados e preservados no Brasil, de modo a ampliar o repertório linguístico e cultural dos estudantes.
Isso pode ser garantido por meio da criação, apreciação e fruição estética nas diversas linguagens artísticas, na perspectiva de ampliar o olhar investigativo e reflexivo no estudante, a partir das vibrações sonoras ( música e ruídos ), das representações criativas ( imaginação e gestualidade ), dos elementos presentes no cotidiano ( imagens estáticas e em movimento ), do movimento rítmico e intencional na expressão corporal das pessoas em sociedade.
E também pela exploração e compreensão do movimento e gestualidade nas práticas corporais, estimulando no estudante a curiosidade intelectual, com autonomia, no âmbito da pesquisa e das práticas corporais na relação inter/intrapessoal.
Cabe destacarmos a importância de o sistema de ensino, por meio das redes estaduais, municipais, federal e distrital, realizar a implementação do Novo Ensino Médio.
Isso, sem dúvida, possibilitará \textbf{uma} prática docente inovadora, tendo como aliados os recursos tecnológicos, de modo a garantir \textbf{uma} aprendizagem significativa e transformadora para os estudantes.
Dessa forma, a área de Linguagens e suas Tecnologias deve considerar, conforme a BNCC ( 2017 ), os cinco campos de atuação dos componentes \textbf{que} integram a mencionada área, a saber : i ) campo da vida pessoal ; ii ) campo das práticas de estudo e pesquisa ; iii ) campo jornalístico-mediático ; iv ) campo de atuação na vida pública ; v ) campo artístico.
É importante destacarmos, ainda, o \textbf{que} é apresentado em cada \textbf{um} dos campos.
O campo da vida pessoal apresenta -se intrínseco a todas as pessoas e propõe > possibilitar \textbf{uma} reflexão sobre as condições \textbf{que} cercam a vida \textbf{contemporânea} e a condição juvenil no Brasil e no mundo e sobre temas e questões \textbf{que} afetam os jovens.
As vivências, experiências, análises críticas e aprendizagens propostas nesse campo podem se constituir como suporte para os processos de construção de identidade e de projetos de vida (... ) \textbf{que} possibilitam \textbf{uma} ampliação de referências e experiências culturais diversas e do conhecimento sobre si ( BNCC, 2017, p. 479 ).
O campo das práticas de estudo e pesquisa insere -se no âmbito intelectual.
Esse campo possibilita ampliar as aprendizagens de cunho científico do estudante, por meio > Da pesquisa, recepção, apreciação, análise, aplicação e produção de discursos/textos expositivos, analíticos e argumentativos, \textbf{que} circulam tanto na esfera escolar como na acadêmica e de pesquisa, assim como no jornalismo de divulgação científica ( BNCC, 2017, p. 480 ).
O campo jornalístico-mediático, \textbf{bastante} presente no cotidiano das pessoas, principalmente dos grandes centros urbanos.
Esse campo se constitui > pela circulação dos discursos/textos da mídia informativa ( impressa, televisiva, radiofônica e digital ) e pelo discurso publicitário.
Sua exploração permite construir \textbf{uma} consciência crítica e seletiva em relação à produção e circulação de informações, \textbf{posicionamentos} e induções ao consumo ( BRASIL, 2017, p. 480 ).
O campo de atuação na vida pública, muito presente e necessário na sociedade, pois > contempla os discursos/textos normativos, legais e jurídicos \textbf{que} regulam a convivência em sociedade, assim como discursos/textos propositivos e reivindicatórios ( petições, manifestos etc. ).
Sua exploração permite aos estudantes refletir e participar na vida pública, pautando -se pela ética ( BRASIL, 2017, p. 480 ).
O campo artístico, \textbf{um} campo democrático das práticas humanas expressas de forma intencional com caráter estético, crítico e reflexivo, sendo considerado > o espaço de circulação das manifestações artísticas em geral, contribuindo para a construção da apreciação estética, significativa para a constituição de identidades, a vivência de processos criativos, o reconhecimento da diversidade e da multiculturalidade e a expressão de sentimentos e emoções.
Possibilita aos estudantes, portanto, reconhecer, valorizar, fruir e produzir tais manifestações, com base em critérios estéticos e no exercício da sensibilidade ( BRASIL, 2017, p. 480 ).
Acrescentamos, ainda, \textbf{que}, segundo a BNCC, a organização por campos de atuação corresponde a três importantes dimensões de formação do \textbf{sujeito} do ensino médio : \textbf{uma} formação estética, \textbf{que} envolve o contato com o \textbf{literário} ; \textbf{uma} formação para o exercício mais direto da cidadania, \textbf{que} envolve a condição de se inteirar dos fatos do mundo, opinar e agir sobre eles ; \textbf{uma} formação \textbf{que} contempla a produção do conhecimento e a pesquisa ( SUGESTÕES CURRICULARES PARA O ENSINO MÉDIO SEDUC–MA[^7 ], 2017 ).
Cabe destacar a importância de observar e desenvolver as atividades acadêmicas, artísticas, sociais e políticas sempre conectadas com os campos de atuação dos componentes.
Atrelado a isso, o trabalho com os itinerários formativos escolhidos pelo estudante ( flexibilização curricular ), garantindo a ele \textbf{uma} formação pautada na escolha, em conformidade com os espaços e tempos escolares oferecidos e na obrigatoriedade das aprendizagens mínimas ( BNCC ).
Acreditamos \textbf{que}, assim, será possível \textbf{uma} educação melhor para todos.
Dessa forma, em conformidade com a BNCC, podemos “ garantir aos estudantes o desenvolvimento de competências específicas ” ( BRASIL, 2017, p. 481 ) \textbf{que} > definem aprendizagens relativas às especificidades e aos saberes \textbf{historicamente} construídos acerca das Línguas, da Educação Física e da Arte ( competências específicas 4, 5 e 6, respectivamente ), \textbf{enquanto} as demais contemplam aprendizagens \textbf{que} \textbf{atravessam} os componentes da área ( BRASIL, 2017, p. 481 ).
Apresentamos, a seguir, a discussão para cada componente curricular da área de Linguagens e suas Tecnologias.
O componente curricular arte, conforme apresentado na BNCC ( 2017 ), tem, em sua estrutura, \textbf{um} instrumento \textbf{norteador} para a \textbf{aplicabilidade} de conteúdos de forma organizada e concisa, contribuindo para “ o desenvolvimento da autonomia criativa e \textbf{expressiva} dos estudantes, por meio da conexão entre racionalidade, sensibilidade, \textbf{intuição} e \textbf{ludicidade} ”.
Além disso, a arte é “ propulsora da ampliação do conhecimento do \textbf{sujeito} relacionado a si, ao outro e ao mundo ” ( p. 474 ).
A arte está presente na história da \textbf{humanidade} desde os tempos mais remotos.
As primeiras pessoas, ao experimentá -la, iniciaram \textbf{um} processo de criação, transformação e descoberta do mundo \textbf{que} as cercava, pois, em suas complexidades ( \textbf{humanidade} e meio circundante ), seguem em constantes inquietações, \textbf{que} se manifestam por meio do fazer criativo.
Isso porque a relação de construção e apropriação do conhecimento, a partir da arte, sempre fez parte de nossas vidas.
Valorizar este componente por sua capacidade de propiciar o desenvolvimento cognitivo em compasso com o emocional é o legado da proposta de ensino da arte por meio da \textbf{abordagem} triangular, defendida no Brasil pela educadora Ana Mae Barbosa.
Esta autora propõe \textbf{que} o ensino de arte no espaço da escola deve ser desenvolvido na perspectiva do “ fazer, fruir e contextualizar ”.
Essa \textbf{abordagem} possibilita as três dimensões necessárias a serem vivenciadas pelo estudante : fazer ( experimentar, criar e recriar ) ; fruir ( apreciar o espetáculo, a música, ler a obra de arte ; para tanto, é \textbf{preciso} conhecer para aprender a gostar ) ; e contextualizar ( analisar, criticar a obra, a partir de conhecimento de mundo, conhecimento estético e reconhecimento de valores culturais ). \textbf{Amparando} a \textbf{aplicabilidade} da arte, a LDB nº 9.394/96 propõe a obrigatoriedade do seu ensino, reconhecendo sua importância no currículo escolar.
De acordo com o artigo 26 e seus parágrafos, > § 2º O ensino da arte, especialmente em suas expressões regionais, constituirá componente curricular obrigatório da educação básica. > > § 6º As artes visuais, a dança, a música e o teatro são as linguagens \textbf{que} constituirão o componente curricular de \textbf{que} trata o § 2º deste artigo ( Redação dada pela Lei nº 13.278, de 2016 ).
Enfatizamos, portanto, a importância do componente curricular arte, considerando, em suas especificidades, as bases para a formação dos/as estudantes : > É na aprendizagem, na pesquisa e no fazer artístico \textbf{que} as \textbf{percepções} e compreensões do mundo se ampliam no âmbito da sensibilidade e se interconectam, em \textbf{uma} perspectiva poética em relação à vida, \textbf{que} permite aos \textbf{sujeitos} estar abertos às \textbf{percepções} e experiências, mediante a capacidade de imaginar e ressignificar os cotidianos e rotinas ( BRASIL, 2017, p. 474 ).
Assim, no ensino médio, devemos compreender \textbf{que} o ensino da arte, sendo \textbf{uma} área de conhecimento \textbf{que} compreende as especificidades das quatro linguagens artísticas ( artes visuais, dança, música e teatro ), deve garantir o diálogo entre si e demais unidades curriculares, no sentido de “ promover o cruzamento de culturas e saberes, possibilitando aos estudantes o acesso e a interação com as distintas manifestações culturais \textbf{populares} presentes na sua comunidade ” ( BRASIL, 2017, p. 474 ).
Entendendo desta forma a arte como \textbf{uma} importante área de conhecimento, necessária à formação integral dos estudantes, como bem evidenciam John Dewey e Paulo Freire, citado por Barbosa, a educação “ é mediatizada pelo mundo em \textbf{que} se \textbf{vive}, formatada pela cultura, \textbf{influenciada} por linguagens, impactada por crenças, classificada pela necessidade, afetada por valores e moderada pela individualidade ” ( BARBOSA, 2010, p. 12 ).
Para Eisner, o trabalho da arte na educação refina os sentidos e alarga a imaginação, potencializando, inclusive, a \textbf{cognição} ( BARBOSA, 2010 ).
Nesse sentido, o ensino da arte, considerando a soma de conhecimentos do universo artístico na história da \textbf{humanidade}, requer \textbf{uma} ampla carga horária no ensino médio, a fim de garantir aos estudantes, no espaço da escola, \textbf{uma} indispensável aprendizagem em arte e seus \textbf{desdobramentos} para \textbf{um} pleno convívio em sociedade.
Para Read ( 2001 ), a arte é a base da educação, contribui para \textbf{um} melhor \textbf{direcionamento} da educação básica, pelo determinante fato de \textbf{que} quando tornamos o estudo prazeroso, obtemos melhores resultados.
Esta satisfação dá -se pelo caráter familiar ao conteúdo estudado, pois é de comum sentido \textbf{que} convivemos com diferentes linguagens da arte desde o nascimento, ela “ (... ) está presente em tudo o \textbf{que} fazemos para satisfazer nossos sentidos ” ( p. 16 ).
Quando o \textbf{autor} define a forma da obra de arte, ele a insere como \textbf{um} fazer essencialmente humano e necessário para o desenvolvimento cognitivo.
O componente curricular educação física, no trabalho com a corporeidade e a motricidade, contribui nos atos de linguagem.
Cabe considerarmos, assim, conforme aponta Sousa Filho ( 2010, p. 1 ), \textbf{que} “ a condição da educação física como componente curricular na história da educação brasileira sempre esteve ligada às tentativas históricas de consolidação de \textbf{um} modelo de política educacional \textbf{ideal} para o país ”.
Com a promulgação da Lei de Diretrizes e Bases da Educação Nacional, tivemos a legitimação da obrigatoriedade da educação física na escola, ou seja, como componente curricular.
Assim, com base no seu art. 26, verificamos \textbf{que} “ a educação física, integrada à proposta pedagógica da escola, é componente curricular obrigatório da educação básica, ajustando -se às faixas etárias e às necessidades da população escolar ” ( BRASIL, 1996 ).
Ressaltamos, portanto, \textbf{que}, de acordo com a LDB nº 9.394/96 e com a promulgação da reforma do ensino médio, por meio da Lei nº 13.415/2017, a educação física está inserida na Base Nacional Comum Curricular ( BNCC ) nos currículos da educação infantil, do ensino fundamental e do ensino médio e deve “ [... ] ser complementada, em cada sistema de ensino e em cada estabelecimento escolar, por \textbf{uma} parte diversificada, \textbf{exigida} pelas características regionais e locais da sociedade, da cultura, da economia e dos educandos ” ( BRASIL, 1996 ).
Assim, na BNCC ( 2017 ), é apresentado \textbf{que}, ao experimentarem práticas da educação física em temáticas como brincadeiras e jogos, esporte, ginástica, dança, lutas e práticas corporais de aventura, os jovens se movimentam com diferentes intencionalidades, construídas em suas experiências pessoais e sociais com a cultura corporal de movimento ( p. 475 ).
Nesse sentido, destacamos \textbf{que} o componente curricular educação física é apresentado na BNCC para o ensino médio como \textbf{uma} \textbf{abordagem} integrada da cultura corporal de movimento na área de Linguagens e suas Tecnologias, aprofunda e amplia o trabalho realizado no ensino fundamental, criando oportunidades para \textbf{que} os estudantes compreendam as \textbf{inter-relações} entre as representações e os saberes vinculados às práticas corporais, em diálogo constante com o patrimônio cultural e as diferentes esferas/campos de atividade humana ( BRASIL, 2017, p. 475 ).
Dessa forma, ressaltamos a importância de tratar temas como o direito ao acesso às práticas corporais pela comunidade, a \textbf{problematização} da relação dessas manifestações com a saúde e o lazer ou a organização autônoma e autoral no envolvimento com a variedade de manifestações da cultura corporal de movimento.
Isso, efetivamente, permitirá aos estudantes a aquisição e/ou o aprimoramento de certas habilidades ( BNCC, 2017 ).
O componente língua portuguesa, conforme previsto na Lei nº 13.415/2017, deve ser oferecido nos três anos do ensino médio.
Assim sendo, as habilidades desse componente também serão organizadas, como no ensino fundamental, por campos de atuação social, sem indicação de seriação.
Essa decisão permite orientar possíveis progressões na definição anual dos currículos e propostas pedagógicas de cada escola ( ORIENTAÇÕES CURRICULARES PARA O ENSINO MÉDIO – MA, 2017 ).
O ensino da língua portuguesa, de acordo com as Diretrizes Curriculares Nacionais ( DCNs ), deve estar voltado para a função social da língua.
Esta é requisito básico para \textbf{que} a pessoa ingresse no mundo letrado, possa construir seu processo de cidadania e, ainda, \textbf{consiga} se integrar à sociedade de forma ativa e o mais autônoma possível ( ORIENTAÇÕES CURRICULARES PARA O ENSINO MÉDIO – MA, 2017 ).
Dessa forma, sobretudo, o componente língua portuguesa possibilitará aos estudantes ( principalmente aos jovens ) experiências significativas \textbf{que} façam sentido às situações de letramentos com práticas de linguagens em diferentes mídias, tendo a escola a incumbência de ampliar situações de significativa aprendizagem.
Na perspectiva de adensamento de conhecimentos iniciados no ensino fundamental, o jovem deve ser preparado para desenvolver maior nível de teorização e análise crítica, principalmente se for levado em consideração o contexto de sua \textbf{territorialidade} e as \textbf{cenas} culturais/sociais em \textbf{que} ele convive.
Nesse aspecto, para ser considerado competente em língua portuguesa, o estudante precisa dominar habilidades \textbf{que} o capacitem a \textbf{viver} em sociedade, atuando de maneira adequada e relevante nas mais diversas situações sociais de comunicação.
Para tanto, precisa saber interagir verbalmente, isto é, precisa ser capaz de compreender e participar de \textbf{um} diálogo ou de \textbf{uma} conversa e produzir textos escritos dos diversos gêneros \textbf{que} circulam socialmente.
Ler e escrever, por suas particularidades formais e \textbf{funcionais}, são também competências mais especificamente desenvolvidas no ambiente escolar.
Tanto os textos escritos de uso mais familiar ( como o bilhete, a carta ) quanto os textos de domínio público ( como o artigo, a notícia, a reportagem, o aviso, o anúncio, o conto, a crônica etc. ) são objeto do estudo sistemático na escola.
Para orientar \textbf{uma} \textbf{abordagem} integrada dessas linguagens e de suas práticas, a área define \textbf{que} os campos de atuação social são \textbf{um} dos seus principais eixos organizadores.
Segundo essa opção, a área propõe \textbf{que} os estudantes possam vivenciar experiências significativas com práticas de linguagem em diferentes mídias ( impressa, digital, analógica ), situadas em campos de atuação social diversos, vinculados com o enriquecimento cultural próprio, as práticas cidadãs, o trabalho e a continuação dos estudos ( BNCC, 2017 ).
O documento curricular do Maranhão tem relação \textbf{dialógica} irrestrita com a Base Nacional Comum Curricular ( BNCC ), fazendo \textbf{uma} ponte direta com as habilidades específicas da língua portuguesa presentes no referido documento, com a cultura e as características específicas do nosso estado.
Dessa forma, cabe \textbf{dizer} \textbf{que} se apresentam as habilidades específicas para língua portuguesa, sugeridas pela BNCC, ofertando possibilidade de diálogo com as demais \textbf{disciplinas}, permitindo \textbf{uma} integração efetiva com os demais componentes da área.
Destaca -se, também, \textbf{que} algumas práticas de linguagem, desenvolvidas em língua portuguesa, dialogam diretamente com todos os componentes, ampliando a construção dessa integração, de modo \textbf{que} não se tenha \textbf{um} componente isolado nem dos campos de atuação, nem das habilidades e competências específicas e, tampouco, fechado em si mesmo.
Dando continuidade à perspectiva investigativa e de abstração adotada no ensino fundamental, a pesquisa e a produção colaborativa precisam ser o modo privilegiado de tratar os conhecimentos e discursos abordados no ensino médio.
Particularmente na área de Linguagens e suas Tecnologias, mais do \textbf{que} \textbf{uma} investigação centrada no desvendamento dos sistemas de signos em si, trata -se de assegurar \textbf{um} conjunto de iniciativas para \textbf{qualificar} as intervenções por meio das práticas de linguagem.
A produção de respostas diversas para o mesmo problema, a relação entre as soluções propostas e a diversidade de contextos e a compreensão dos valores éticos e estéticos \textbf{que} \textbf{permeiam} essas decisões devem se tornar foco das atividades pedagógicas. ( BNCC, 2017 ).
Para isso, é fundamental \textbf{que} sejam garantidas, aos estudantes, oportunidades de experienciar fazeres cada vez mais próximos das práticas da vida acadêmica, profissional, pública, cultural e pessoal, e situações \textbf{que} \textbf{demandem} a articulação de conhecimentos, o planejamento de ações, a auto-organização e a negociação em relação a metas.
Tais oportunidades também devem ser orientadas para a criação, para o encontro com o inusitado, com vistas a ampliar os \textbf{horizontes} éticos e estéticos dos estudantes.
Destacamos \textbf{que} é importante \textbf{que} os jovens, ao explorarem as possibilidades \textbf{expressivas} das diversas linguagens, possam realizar reflexões \textbf{que} envolvam o exercício de análise de elementos discursivos, composicionais e formais de enunciados nas diferentes semioses – visuais ( imagens estáticas e em movimento ), sonoras ( música, ruídos, sonoridades ), verbais ( oral ou \textbf{visual-motora}, como Libras e escrita ) e corporais ( gestuais, cênicas, dança ).
Afinal, muito por efeito das novas tecnologias da informação e da comunicação ( TIC ), os textos e discursos \textbf{atuais} organizam -se de maneira híbrida e multissemiótica, incorporando diferentes sistemas de signos em sua constituição ( BRASIL, 2017 ).
O componente curricular língua inglesa, conforme o parágrafo 4º do artigo 35 -A da LDB, também é obrigatório no ensino médio porque essa língua deve ser compreendida como sendo de uso mundial, pela multiplicidade e variedade de usos, usuários e funções na \textbf{contemporaneidade}, assim como definido na BNCC do Ensino Fundamental – Anos Finais.
Com a reforma do ensino médio, a língua inglesa ganha destaque, \textbf{uma} vez \textbf{que} a Lei nº 13.415/17 incluiu o art. 35 -A na LDB nº 9.394/96, determinando \textbf{que} > Os currículos do ensino médio incluirão, obrigatoriamente, o estudo da língua inglesa e poderão ofertar outras línguas estrangeiras, em caráter optativo, preferencialmente o espanhol, de acordo com a disponibilidade de oferta, locais e horários definidos pelos sistemas de ensino.
Cabe destacarmos \textbf{que}, no Brasil, as \textbf{abordagens} de ensino da língua inglesa e sua estrutura curricular sofreram constantes mudanças, em consequência da organização social, cultural, política e econômica ao longo do tempo.
As propostas curriculares e os \textbf{métodos} de ensino vêm sendo redefinidos para atenderem às expectativas e demandas sociais da atualidade, de forma a assegurar a aprendizagem ( ORIENTAÇÕES CURRICULARES PARA O ENSINO MÉDIO – MA, 2017 ).
De acordo com a Base Nacional Comum Curricular, a língua inglesa é \textbf{vista} como > língua franca, \textbf{uma} língua de comunicação internacional utilizada por falantes espalhados no mundo inteiro, com diferentes repertórios linguísticos e culturais.
Essa perspectiva permite questionar a visão de \textbf{que} o único inglês correto – e a ser ensinado – é aquele falado por estadunidenses ou britânicos, por exemplo.
Desse modo, o tratamento do inglês como língua franca o desvincula da noção de pertencimento a \textbf{um} determinado território e, consequentemente, a culturas típicas de comunidades específicas ( BRASIL, 2017, p. 198 ).
Fica evidente, portanto, a preocupação da BNCC quanto ao desenvolvimento do ensino da língua inglesa em consonância com seu caráter universal, o \textbf{que} ganha mais relevância em \textbf{um} contexto social marcado pela globalização e pela competitividade.
Segundo a BNCC ( BRASIL, 2017, p. 201 ), estudar a língua inglesa permite ao estudante “ ampliar \textbf{horizontes} de comunicação e de intercâmbio cultural, científico e acadêmico e, nesse sentido, abre novos percursos de acesso, construção de conhecimentos e participação social ”.
Portanto, a aquisição de \textbf{uma} língua estrangeira propicia aos jovens conhecer não apenas outras culturas, como também ser capaz de compreender a sua própria e de ser agente modificador do seu entorno sociocultural.
Destacamos, também, \textbf{que} a aprendizagem da língua inglesa na \textbf{contemporaneidade} \textbf{configura} -se na preparação do estudante para ser \textbf{um} cidadão do mundo, pertencente a \textbf{um} território sem tantas \textbf{fronteiras} linguísticas, em \textbf{que}, por meio da aquisição de \textbf{uma} língua estrangeira ( inglês ), ele poderá estar inserido num mundo cada vez mais \textbf{globalizado} e diverso.
Dessa forma, o Documento Curricular do território maranhense ( DCTM ) comunga da mesma proposição colocada na Base Nacional Comum Curricular ( BNCC ), trabalhando com as competências da área de Linguagem e suas tecnologias, bem como as habilidades \textbf{que} são específicas da língua inglesa, respeitando as idiossincrasias presentes em nosso estado, além de propor \textbf{uma} interação interdisciplinar e transdisciplinar dentro da área a \textbf{que} pertence e com as outras áreas do conhecimento.
Reforçamos, também, o fato de \textbf{que} o estudo da língua inglesa se faz presente em \textbf{um} diálogo constante com todos os componentes curriculares, construindo, assim, \textbf{uma} interação \textbf{que} perpassa pela quebra do isolamento disciplinar.
Quanto aos Temas \textbf{Contemporâneos} também presentes na BNCC, o ensino da língua inglesa no território maranhense deve -se constituir na perspectiva de \textbf{um} aprendizado voltado para as práticas sociais, para \textbf{que} a língua estrangeira seja \textbf{uma} ferramenta para ampliar cada vez mais o repertório social e cultural no nosso estudante com vistas a promover na escola \textbf{um} espaço de discussão gradativamente mais democrático e plural.
Formando, assim, \textbf{uma} juventude mais cidadã, \textbf{identitária} e \textbf{comprometida} com as novas transformações de \textbf{uma} sociedade \textbf{que} \textbf{exige} dos nossos estudantes \textbf{uma} capacidade empreendedora e de soluções mais criativas.
Sendo assim, podemos \textbf{dizer} \textbf{que} as aprendizagens em língua inglesa permitirão aos estudantes do estado do Maranhão usar essa língua para aprofundar a compreensão sobre o mundo em \textbf{que} \textbf{vivem}, explorar novas perspectivas de pesquisa e obtenção de informações, expor ideias e valores, \textbf{argumentar}, lidar com conflitos de opinião com crítica e com criatividade para solucionar problemas, entre outras ações relacionadas ao seu desenvolvimento cognitivo, linguístico, cultural e social.
Desse modo, o ensino de língua inglesa os ajuda a ampliar sua capacidade discursiva e de reflexão em diferentes áreas do conhecimento ( BRASIL, 2017 ).
Quanto à língua espanhola, no Brasil, sua inserção na educação básica, do ponto de vista do percurso histórico, não é fato recente, remontando à década de 1940, com a chamada Reforma Capanema, sendo, no entanto, retirada do currículo duas décadas depois, nos anos 60.
No contexto sul-americano, a partir dos anos 90, o fortalecimento do Mercado Comum do Sul ( Mercosul ) proporcionou ampla difusão do espanhol, fazendo com \textbf{que} cidadãos comuns, instituições de ensino, empresas públicas e privadas procurassem romper \textbf{fronteiras} linguístico-histórico-culturais e demonstrassem interesse pelo estudo desse idioma.
Seguindo a trajetória histórica do processo de inclusão do espanhol na educação básica brasileira, em 1996, a LDB abriu espaço, outra vez, para o ensino desse idioma no âmbito da educação básica, ao determinar, em seu artigo 26, parágrafo 5º, \textbf{que} na parte diversificada do currículo deveria ser incluído, obrigatoriamente, a partir da 5ª série, o ensino de pelo menos \textbf{uma} língua estrangeira \textbf{moderna}, ficando a escolha a cargo da comunidade escolar ( BRASIL, 1996 ).
Destaca -se, também, a Lei nº 13.415 de 16 de fevereiro de 2017 ( art. 3º ), Resolução CEE/MA 277/2021 ( art. 7º ), Resolução nº 3, de 21 novembro de 2018 ( art. 11 ), estabelecem a oferta da língua espanhola, em caráter optativo, de acordo com a disponibilidade de oferta, locais e horários das instituições de ensino.
Assim sendo, entendemos \textbf{que} o trabalho com o ensino do espanhol pode perpassar por diferentes possibilidades de aprendizagens e ampliação de conhecimentos, do ponto de vista de sua história e até mesmo considerando a necessidade de compreensão da cultura dos países hispanohablantes. \textbf{Ademais}, tem -se a questão das diversidades linguísticas e variações da língua \textbf{que}, por trás de tudo isso, há culturas, costumes e gostos distintos \textbf{que} não devem ser desconsiderados.
Para isso, é \textbf{preciso} considerar \textbf{que} a prática de ensino nas escolas seja realizada de forma a promover \textbf{uma} educação \textbf{que} leve em consideração a natureza reflexiva, dinâmica, política e social dos \textbf{atores} envolvidos, em sintonia com as estruturas sociais vigentes, como forma de propiciar aprendizagens significativas para o estudante.
No \textbf{que} tange ao ensino da língua espanhola no estado do Maranhão, torna -se importante a compreensão da relevância deste componente nas áreas econômicas, sociais e culturais desta unidade da federação.
Nessa linha de pensamento, evidenciam -se, primeiramente, as novas perspectivas \textbf{que} o mercado de trabalho demanda, sendo, dessa maneira, o aprendizado do idioma espanhol de grande importância para o fortalecimento da qualificação para o mercado laboral.
Nesse contexto, o avanço econômico experimentado pelo estado demonstra a necessidade do aprendizado desse idioma com vistas à ampliação das práticas linguísticas e sociais dos estudantes, sobretudo em \textbf{face} das potencialidades turísticas maranhenses e do contexto de globalização vigente, \textbf{que} estimula os intercâmbios comerciais e culturais e, consequentemente, novos contatos linguísticos.
Apresentamos \textbf{que} a área de Linguagens e suas Tecnologias está organizada com considerações gerais acerca dos componentes curriculares \textbf{que} a formam, bem como com as competências e as habilidades específicas de cada componente.
Destacamos \textbf{que}, nessa área, temos sete competências, \textbf{que} apresentam suas respectivas habilidades específicas.
Portanto, dentro da arquitetura curricular, no \textbf{que} tange à área de Linguagens e suas Tecnologias, os conteúdos estarão \textbf{embasados} de forma a contribuir para o alcance dos objetivos dos estudantes : continuidade dos estudos, atuação no mundo do trabalho, resolução de demandas complexas da vida cotidiana e exercício pleno da cidadania.
Nessa perspectiva de ensino integrado, conectado e de caráter democrático, por meio da oferta e das escolhas dos estudantes, podemos apresentar \textbf{uma} nova escola, com \textbf{um} novo ensino \textbf{que} ganha mais sentido nas vidas dos estudantes brasileiros, especialmente do território maranhense. [ ^7 ] : Trata -se de \textbf{um} Caderno de Orientações Curriculares para o Ensino Médio, produzido pela SEDUC -MA, \textbf{que} faz parte do Plano Mais IDEB.

% matched lemmas: abordagem, ademais, amparar, ancorar, aplicabilidade, argumentar, ator, atravessar, atual, autor, bastante, cena, cognição, comprometido, configurar, conseguir, contemporaneidade, contemporâneo, demandar, desdobramento, dialógico, direcionamento, disciplina, dizer, embasar, enquanto, exigir, explicitar, expressivo, face, fronteira, funcional, globalizado, historicamente, horizonte, humanidade, ideal, identitária, imbricar, influenciar, inter-relação, intuição, literário, ludicidade, moderno, método, norteador, percepção, permear, popular, posicionamento, preciso, preconizar, problematização, qualificar, que, ribeirinho, sujeito, territorialidade, transdisciplinaridade, um, ver, visual-motor, viver
\end{textsample}
